\section{Discussion and Limitations}
\label{sec:discuss}

\subsection{Summary of Findings Across Three Studies}
\label{subsec:discuss_summary}

The AI-READI port produces a viable participant-held-out streaming forecaster
and improves over persistence on matched anchors.
The improvement over persistence is real but not large, which is the expected
limitation of residual-current CGM forecasting: current glucose already carries
most of the short-horizon predictive signal.

\textbf{Study 1} (Section~\ref{sec:study1_cf_audit}) establishes that AI-READI
cannot support a meal counterfactual with any constructible label.
The glucose-defined arm is sign-inverted ($\Delta\theta_{60} \approx -35$~mg/dL
versus observed $+22.2$~mg/dL) because the label selects precisely the anchors
where the forecast-only path already anticipates a rise.
The wearable arm is a near-zero null ($\Delta\theta_{60} \approx -0.21$~mg/dL),
and a dose-matched sweep to model-space 6.25 yields $-0.77$~mg/dL.
The three convergent findings are not independent causal proofs, but they rule out
the main alternative explanations (dead channels, dose magnitude, encoding artefacts).
The required next step is exogenous meal annotation from a dataset with prospective
logging (T1DEXI or CGMacros).

\textbf{Study 2} (Section~\ref{sec:study2}) takes a different approach: rather
than attempting a meal counterfactual, it asks whether deployable model-surprise
signals can correct forecasts \emph{after} the base model's error becomes
observable.
The rich correction head (Ridge full, split up/down; 19 causal trailing features)
beats the frozen base by $0.175$~mg/dL on test triggered windows
(95\% CI $[-0.259,\;-0.094]$, CI excludes zero) and beats the best
non-learning baseline (decaying shift) by $0.172$~mg/dL
(CI $[-0.238,\;-0.109]$).
The practical contribution is an interpretable, online correction mechanism:
it fires on approximately $9.3\%$ of deployment anchors (triggered rate at $\tau=0.2$),
requires no meal labels or future covariates, and can be retrained without
touching the frozen base.
The improvement is modest in absolute terms.
Direction-specific heads (separate up/down) outperform a combined head; Ridge
and HGBT are statistically tied; the richer causal feature set beats a simple
$z_t$-only correction but the margin is interpretable, not surprising.

\textbf{Study 3} (future work) asks whether the scenario-aware channel supports
a causal meal counterfactual on a dataset with prospective meal logging such as
T1DEXI or CGMacros.
AI-READI does not serve this role: no constructible label carries pre-rise
causal information orthogonal to the CGM slope.

\subsection{Interpretability}
\label{subsec:discuss_interp}

Section~\ref{sec:interp} reports two complementary views of frozen-base model
reliance on the full held-out test split.
At the output level, recent glucose history is the dominant driver
(36.7\% of LOMO-attributed influence), with secondary contributions from static
clinical metadata, site/cohort, time/calendar, data-quality indicators, and
heart-rate/physiology.
Data quality and site/cohort together account for 17.4\% of output-level
influence, which is predictive but also cautionary: part of the model's signal
comes from collection structure rather than purely physiological variation.
The meal-proxy diagnostic channel registers 0.0\% output-level influence,
consistent with Study~1 and arguing against any AI-READI meal-counterfactual
interpretation.

The hidden-state view gives a different but compatible picture.
Among the six groups that build the recurrent Mamba state $h_t$, glucose
history, heart-rate/physiology, and data-quality/missingness carry comparable
state-construction weight, while sleep, wearable activity/exercise, and stress
also contribute.
Thus the model's internal state is richer than a glucose-only trace, even
though the final forecast remains most sensitive to glucose history.

\subsection{Limitations}
\label{subsec:discuss_limitations}

\paragraph{Dataset.}
AI-READI lacks timed insulin doses, insulin-on-board records, and carbohydrate
quantity logs.
The variable \texttt{med\_insulin} is static therapy metadata, not an action
covariate.
The predicted meal flag is noisy and potentially retrospective.
No confirmed meal timestamps are available; the event-enriched proxy used in
Study~2 is a glucose-level proxy and is not distribution-matched to its control
group.

\paragraph{Forecasting target.}
The residual-current target can hide weak dynamic learning if persistence is not
separately reported.
The model does beat persistence, but the margin is modest at every horizon
(Section~\ref{sec:results}).

\paragraph{Personalisation.}
Matched-anchor personalisation is flat in the current artifacts; offset correction
worsens MAE.
Participant-level averages are necessary because anchor-weighted metrics
overrepresent participants with longer clean segments.

\paragraph{Study 2 scope.}
The correction head is evaluated on triggered windows only ($\approx 9.3\%$ of
deployment anchors).
The event-vs-control comparison in Section~\ref{subsec:study2_event} is
descriptive and structurally confounded: the event definition selects on glucose
level, so the two groups differ by construction.
The correction head is slightly harmful for extreme downward residuals
($r_t < -20$~mg/dL, 77 test anchors) and for narrow prediction-interval windows
(Q1); these cases should be handled with explicit abstention in deployment.

\paragraph{Interpretability scope.}
The output-level modality attribution is based on 153{,}591 dense anchors
across 381 test streams (221 participants), the full held-out test split.
The $h_t$ joint-group norm applies only to the six modality groups that enter
the streaming encoder state; static/context and direct-decoder inputs are
interpreted through output-level LOMO instead.
Both analyses are model-reliance audits, not causal effect estimates.

\paragraph{Site and phenotype effects.}
Site and phenotype subgroup effects require follow-up audits before being
interpreted as biological effects rather than cohort composition or domain shift.
The insulin-dependent subgroup in Study~2 failure-mode analysis comprises only
8 participants (1{,}498 anchors); results in that cell should not be
generalised.

\subsection{What a True Meal Counterfactual Requires}
\label{subsec:discuss_true_counterfactual}

A scientifically valid ``what if I eat now?'' counterfactual requires exogenous
meal information not derivable from the CGM signal: prospective meal logs,
carbohydrate quantities, timed insulin boluses, or a standardised meal-challenge
dataset.
Study~1 demonstrates that approximating this signal from glucose-shape or
wearable-activity features fails sign agreement with the observed natural
experiment.
Study~2 provides a deployable alternative that avoids this requirement entirely:
it corrects forecasts using only causal residual and shape features available
at trigger time, without claiming to model the meal causal mechanism.
Future work connecting Study~2 event-aware triggers to prospective meal logs
(e.g., from CGMacros) would enable a direct comparison of the two approaches
on the same participants.
